% =====================================================================
% Analysis section, four research questions.
%
% Structure follows the PriVE-Bench pattern: a bold run-in question, two or
% three paragraphs of evidence, a closing sentence stating what the RQ
% established. RQ1 establishes that the failure is real and diagnosable,
% RQ2 tests the simplest intervention, RQ3 tests whether it transfers to a
% second reasoning interface, RQ4 decomposes the averages.
%
% \pending{} marks numbers that depend on the answerability cleaning pass
% and MUST be refreshed before submission. See the note above RQ1.
% =====================================================================
\newcommand{\pending}[1]{#1}   % 定稿前改成 \textcolor{red}{#1} 自查，然后删掉

\section{Analysis: From Static Reasoning to Structured and Agentic Methods}

\paragraph{Experimental setup.}
We evaluate seventeen frozen systems on all $10{,}198$ items at temperature $0$ under a fixed
seed. Responses from which no option letter can be recovered are scored wrong rather than
dropped, so the denominator is the full benchmark for every system. CFA items carry three
options and FRM items four, so raw accuracy is not comparable across tracks and we report
chance-normalised accuracy alongside it,
\begin{equation}
\mathrm{Acc}^{*}=\frac{\mathrm{Acc}-c}{1-c},
\qquad c=\tfrac{1}{3}\ (\text{CFA}),\quad c=\tfrac{1}{4}\ (\text{FRM}).
\end{equation}
RQ2 and RQ3 hold the answering system fixed and change only the evidence it is given, so every
comparison is paired at the item level. Against each family baseline we report
\begin{equation}
\begin{aligned}
\text{Rescue} &= \#(\text{base wrong},\ \text{method right}),\\
\text{Hurt}   &= \#(\text{base right},\ \text{method wrong}),\\
\text{Net}    &= \text{Rescue}-\text{Hurt},\\
\text{Persist}&= \#(\text{base wrong},\ \text{method wrong}),
\end{aligned}
\end{equation}
tested with an exact McNemar test under Benjamini-Hochberg correction, with item bootstrap
intervals in Appendix~\ref{app:bootstrap}.

\paragraph{RQ1: What makes \textsc{FinExam-10K} challenging for current models?}
Not the difficulty the exams already publish. We assign each item an empirical band from the
consensus of the seventeen systems, splitting them into the ten proprietary and the seven
open-weight systems and weighting the two groups equally so that band membership is not decided
by whichever group is larger. This yields $6{,}578$ easy, $2{,}183$ medium and $1{,}437$ hard
items. Cram\'er's $V$ between this axis and the exams' own level labels is $0.174$, so the two
orderings are largely independent and a benchmark that reports only CFA and FRM levels is not
reporting difficulty in the sense that separates current systems. Split-half reliability of the
consensus score is $0.805$ Spearman and $0.892$ after the Spearman-Brown correction, and the
worst leave-one-model-out stability is $0.983$, so the axis is a property of the items. Because
the bands are defined by model consensus, all band-conditioned results in the main text are
scored leave-one-out, and Appendix~\ref{app:circularity} repeats the analysis with DeepSeek-R1
and GPT-4o removed from the panel entirely, since those two systems also appear as the
backbones in RQ2 and RQ3.

The bands separate systems where the exam stages do not, and we report the hard band twice
because a quarter of the corpus cannot be answered from its own text (Section~\ref{sec:quality}).
Because CFA items carry three options and FRM items four, a system's hard-band score under
random responding follows a Poisson-binomial rather than a binomial law, and we evaluate that
null exactly and correct across the seventeen systems with Benjamini-Hochberg. Over the full
$10{,}198$ items, scored leave-one-out, $16$ of $17$ systems fall below chance and $15$ do so
significantly at $q<0.05$. The band definition is unchanged throughout: the consensus score is a
per-item quantity, so restricting attention to a subset never moves an item between bands. What
changes is which Hard items we score on. Of the $1{,}437$ Hard items, $372$ can be answered from
their own text and $1{,}065$ cannot (Section~\ref{sec:quality}). Scored on that $372$-item
subset the picture separates: $11$ systems remain significantly below a chance rate of $30.5$,
while three now sit significantly above it, led by Gemini-3.1-Pro at $60.3$. The full-band figure
overstates the failure, because no system can answer an item whose data was never delivered, and
the context-complete figure is the one to cite.

The cleaning strengthens rather than weakens the claim that these failures are shared, but the
comparison requires the right null. A uniform floor of $0.500$ is the asymptotic expectation, not
the finite-sample one: when only a handful of systems err on a three-option item, the modal share
exceeds one half by chance alone. We therefore compute each item's own expected modal share under
uniform allocation of its observed errors across its own distractors, by exact enumeration. The
median expectation is $0.623$. Observed concentration on the answerable hard band has a median of
$0.920$, a paired median excess of $+0.326$, and $348$ of $372$ items, $93.5$ percent with a
$95\%$ interval of $[90.6, 95.6]$, exceed their own null. A Wilcoxon signed-rank test on the
paired differences gives $z=16.2$ at $p<10^{-15}$. On items that can be solved, systems are not
spreading their mass over the distractors, they are being pulled to the same wrong answer.

Where the residual difficulty sits depends entirely on doing this correction, and we report the
uncorrected answer as a cautionary case. Over the full corpus CFA Level II holds $39.2$ percent
of the hard band against $21.7$ percent of the benchmark, which invites the conclusion that
multi-step valuation and case reasoning is the bottleneck. That conclusion is an artefact.
Level II is the stage whose case vignettes were most often lost, $54.6$ percent of its items are
unanswerable, and on the answerable subset it supplies only $15.9$ percent of hard items, below
its share of the corpus. The corrected hard band is led by CFA Level I at $36.6$ percent, with
FRM Part II and Part I next at $17.7$ and $17.5$ percent. A benchmark paper that skipped the
answerability audit would have reported a bottleneck that does not exist.

The other asymmetry survives correction. Every finance-specialized system is beaten by a smaller
general-purpose open model: \texttt{gpt-oss-20b} reaches $70.06$ overall against $68.07$ for the
strongest finance-tuned system. Part of that gap is not reasoning. Finance-tuned systems produce
the overwhelming majority of unparsable responses, $1{,}224$ for \texttt{Fin-R1-7B} and $825$ for
\texttt{ODA-Fin-RL-8B} against at most $231$ for any general-purpose model, so on this evidence
domain fine-tuning costs instruction-following without buying domain accuracy. Thus RQ1 shows
that the benchmark isolates a difficulty axis the published exam levels do not predict, and that
the hard band is a regime where systems fail together rather than independently.

\paragraph{RQ2: How far can structured retrieval and verification improve DeepSeek-R1
chain-of-thought?}
Barely at all, and the null survives every configuration we tried. Using DeepSeek-R1 we compare
direct chain-of-thought against four retrieval configurations that vary the retriever, the query
formulation and the selection policy, each over the identical $10{,}198$ item IDs
(Table~\ref{tab:r1-variants}). Accuracy spans $0.62$ points, from $79.22$ to $79.84$ against a
$79.75$ baseline, or $69.78$ to $70.68$ against $70.55$ in $\mathrm{Acc}^{*}$, and no
configuration differs significantly from the baseline after correction.

The aggregate null is a cancellation, not an absence of effect. Every configuration moves about
a thousand items. The best of the four, a graph reranker trained on FinanceReasoning's own
function relevance labels, rescues $509$ items and breaks $500$ for a net of nine. The published
heuristic ordering rescues $501$ and breaks $535$. Removing the judge and injecting the top ten
graph neighbours directly rescues $496$ and breaks $550$, the worst of the set, which at least
confirms the judge is doing something. Stratifying by band shows the same cancellation inside
every band rather than a hard-band gain masked by an easy-band loss. Thus RQ2 shows that under
chain-of-thought the retrieval stack relocates a large amount of probability mass without moving
accuracy, and that learning the ranking from the source benchmark's own supervision does not
change that.

\paragraph{RQ3: Do the same interventions transfer to GPT-4o program-of-thought?}
No. The identical retrieval stack that was neutral under chain-of-thought is significantly
harmful under program-of-thought. Both routes fall below the $69.37$ direct baseline, function
retrieval to $68.09$ and the FunctionGraph-RAG route to $68.06$, at $q=0.0001$ and $q=0.0002$, and
the two routes are indistinguishable from each other at $q=0.954$
(Table~\ref{tab:gpt4o-variants}). Because DeepSeek-R1 with chain-of-thought and GPT-4o with
program-of-thought differ in both the backbone and the reasoning protocol, we do not attribute
the difference to the protocol alone. The defensible statement is that structured augmentation
exhibits different rescue and failure profiles under the two interfaces, which suggests that
retrieval effectiveness depends on how the backbone represents and executes intermediate
reasoning.

The program-of-thought interface exposes where the loss is incurred, and it is not in retrieval.
Injecting any function raises the rate at which the restricted execution VM refuses the generated
program from $2.4$ percent to about ten percent, and the rate is flat in the number of functions
injected, so it is the presence of shown code rather than its quantity that does the damage
(Table~\ref{tab:pot-mechanism}). $273$ items that executed cleanly without retrieval are refused
with it, and $144$ of those had been answered correctly, which alone exceeds the entire net loss
of $132$ items. The dominant new refusal is a nested function definition, $130$ of the $273$:
shown function definitions as context, the model begins writing nested definitions, which the
sandbox forbids. The graph route halves this specific harm, from $197$ refusals to $100$, and
recovers no accuracy, so a second channel is operating in which the injected functions distract
the reasoning itself.

Retrieval reach caps the exercise independently of these effects. Contriever recall at $30$ on
this corpus is $0.626$ under the best query formulation we measured, and the judge selects no
function at all for $7{,}002$ items, $68.7$ percent of the benchmark, so for roughly a third of
the items the gold function is never retrieved and no downstream selector can recover it. Thus
RQ3 shows that a retrieve-then-inject pipeline can be net negative even when each component
behaves as designed, because the cost is paid in the generation stage rather than the retrieval
stage.

\paragraph{RQ4: Where does the gain from structured retrieval actually sit?}
Not in any single configuration, and not in an average over them. Across the five
chain-of-thought conditions no variant beats the baseline by more than $0.09$ points, yet an oracle
that could pick the best condition per item reaches $89.31$ against the $79.75$ baseline, a headroom of
$9.56$ points. The five program-of-thought conditions behave the same way, with a per-item oracle at
$76.83$ against $69.37$, a headroom of $7.46$ points. Taking both families together the oracle
reaches $92.59$. The variants are therefore not redundant and the stack is not saturated. The
information needed to answer these items is already being produced by methods we have run. What
no static rule extracts is the decision of which condition to trust on which item, and that decision
cannot be made before the item is seen.

The verifier ablation shows what unlocking that decision requires. We fire a conditional verifier
on the $1{,}392$ items where the two retrieval routes disagree. Given only the two answer
letters, $50$ input tokens, it abstains on every conflict and changes nothing, which is the only
defensible reply to that prompt. Given the question, the options, both generated programs, both
execution results and the injected function identifiers, $691$ input tokens, the same model at
the same temperature and seed changes $434$ answers, rescuing $193$ and breaking $109$ for a net
gain of $84$ at $p=1.5\times10^{-6}$. It recovers $65$ percent of the retrieval penalty, after
which the residual gap to the no-retrieval baseline is no longer significant at $p=0.158$. Its
error taxonomy points back at RQ3's mechanism rather than at a knowledge gap:
\texttt{unsupported\_assumption} at $560$ and \texttt{formula\_selection\_error} at $167$ both
describe a program that ran and returned a number derived from the wrong premise.

What remains after every static method is small and sharply concentrated. Define
\begin{equation}
W_{\mathrm{persistent}}=\bigcap_{m\in\mathcal{M}}\{i:\hat y_{i,m}\neq y_i\}
\end{equation}
over all ten conditions in both families. Over the full corpus it contains $756$ items, but $65$ percent
of those are simply unanswerable, so the meaningful figure is the $261$ that survive the
answerability filter, $3.4$ percent of the answerable corpus. Of these $203$ are hard, which is
$54.6$ percent of the entire answerable hard band, and exactly one is easy. Enrichment is highest
in CFA Level III at $2.16\times$ and FRM Part II at $1.94\times$, the two stages that ask for
judgment under composed constraints rather than for a single valuation. We audit a stratified
sample of $261$ items. The audit categories are ordered so that data defects are excluded before
reasoning failures are attributed: an item is first checked for
whether it is answerable as shipped, and only then classified as retrieval miss, judge or
selector miss, retrieved evidence present but unused, multi-step planning failure, rule and
exception composition failure, numerical execution failure, verifier miscorrection, or formatting
failure. Thus RQ4 shows that the binding constraint is not coverage of the function library but
the absence of any mechanism that inspects what the earlier stages produced and chooses between
them.

\paragraph{From static augmentation to agentic reasoning.}
RQ2 and RQ3 measure matched interventions under two reasoning interfaces and find neutrality in
one and harm in the other. RQ4 then locates the gain that neither reports: it sits between the
conditions rather than inside any one of them, and reaching it requires a per-item decision that a
fixed pipeline has no place to make. The informed verifier is the smallest possible instance of
that decision, a single inspection step that recovers most of the retrieval penalty, and its
gap to the oracle is the space an adaptive policy would occupy.

This motivates an agentic formulation in which the model chooses among answering directly,
retrieving function or graph evidence, executing code, invoking a verifier, revising and
abstaining, rather than receiving a context assembled before the problem is seen. The two
benchmarks compose for that study without overlapping. FinanceReasoning supplies $2{,}238$
open-ended numeric problems with executable reference solutions and direct \texttt{function\_id}
relevance labels, of which $511$ are usable under our candidate protocol, which is enough
supervision to fit a policy over lookup and verification actions. We emphasise that this is
relevance-label supervision, not trajectory supervision, since the source dataset does not
publish reasoning traces. \textsc{FinExam-10K} then supplies $10{,}198$ held-out multiple-choice
items in a different answer format from a disjoint item pool, with an empirical difficulty axis
and a hard band no current system clears. A policy frozen after training on the former and
evaluated once on the latter should be reported not only on accuracy but on tool calls, tokens,
latency, rescue, hurt, and regret against the per-item oracle established in RQ4. After release
the $5{,}110$ public items serve as the development split and the $5{,}088$ held-out items as the
agentic leaderboard test set.
